\section{Introduction}
The psychological significance of a musical artefact cannot be inferred solely from its acoustic features. A person may choose a song to sustain a mood, communicate an experience, redirect attention or reconsider a situation. These purposes can coexist, and the same musical activity may have different meanings across occasions. Research on music and mood regulation distinguishes functions of engagement \citep{saarikallio2007,baltazar2016,baltazar2019}, whereas accounts of music-evoked emotion describe routes through which a listener may experience an affective response \citep{juslin2008}. Experiencing an emotion, expressing an emotion and intentionally regulating an emotion are related but distinct processes. A creator may regard a sad song as an accurate representation of an experience without seeking sadness or relief.

Text-conditioned music generation extends the range of available creative operations \citep{copet2023}. Users can describe desired content or style, choose among candidate outputs and refine an artefact without performing every compositional operation themselves. The AI Song Contest documented the iterative exploration and curation involved in human--AI songwriting \citep{huang2020}. Qualitative research has examined how creators negotiate authorial positioning and creative agency when using such systems \citep{zhao2026}; recent experimental work has also examined relationships between automation, task load, ownership and state agency \citep{he2026}. These studies make broad claims that creative agency in AI music is a new research topic untenable. The present study instead focuses on the connection between a concrete personally meaningful episode, specific creative decisions and the emotional significance attributed to its outcome.

\subsection{Experience, expression and personal meaning}
Emotion regulation may involve changes to situations, attention, appraisals and responses \citep{gross2015}. Writing a description, transforming its perspective and listening to the resulting song may each be meaningful to a creator. Expressive writing research provides an adjacent paradigm \citep{pennebaker1986}, while research on self-distancing distinguishes potentially constructive reflection from immersion in negative experiences \citep{kross2011}. Neither establishes that AI-generated songs produce the same consequences. Likewise, findings on autobiographical salience during music listening \citep{janata2009} motivate questions about newly produced autobiographical artefacts without demonstrating that creating such artefacts changes memory or identity.

The proposed AI-assisted Music for Emotion Regulation (AIM-ER) framework distinguishes a regulatory goal, creative operations, proximal psychological mechanisms and outcomes. Here these distinctions guide attention to possible relationships without imposing a fixed sequence or presupposing emotion regulation as the user's motive. A person may begin with curiosity, recognise emotional relevance only after hearing an output and then revise it; another may enter with a specific experience and abandon an output that fails to represent it. Participants' own terms and interpretations take priority over theoretical labels. A retrospective claim of relief will be described as perceived relief rather than an experimentally established change.

\subsection{Agency, effort and ambiguous outcomes}
Agency concerns perceived influence over one's actions; psychological ownership concerns the extent to which the artefact feels one's own. Both differ from attributed contribution and legal authorship. Self-determination theory offers a related account of autonomy and competence \citep{ryan2000}, but does not imply that more available controls always increase either construct. A creator might welcome automation, feel constrained by inconsistent output, or value surprise as part of the creative process. We will attend to both the work performed by the creator and the meaning assigned to that work, including delegation, failed attempts, selection, revision and abandonment.

Repeated focus on a difficult experience may support reflection or accompany rumination \citep{nolen2008}. Evidence on diverse music-use patterns cautions against treating music as uniformly adaptive or maladaptive \citep{hums2015,carlson2015}. Regulatory flexibility also emphasises the relevance of context and adjustment to feedback \citep{bonanno2013}. Although musical activities can have regulatory benefits in some settings \citep{peters2024}, such evidence does not establish a benefit specific to AI-assisted generation. Interviews will therefore allow for unpleasant, mixed, unchanged and unintended experiences, not only relief or successful expression.

\section{Aim and research questions}
The study aims to examine how adults describe the translation of personal experience into AI-assisted musical artefacts and the significance they attribute to this process. It is an exploratory inquiry, not a confirmatory test of AIM-ER propositions.
\begin{enumerate}
\item[RQ1.] What motivates creators to initiate a particular session, and how do their goals change during it?
\item[RQ2.] How do creators translate experiences into prompts, lyrics, stylistic constraints, selection and revision, and how do they interpret mismatches between intended and generated content?
\item[RQ3.] How do they describe their contribution, perceived control, ownership and the personal resonance of an artefact, including instances in which these diverge?
\item[RQ4.] What emotional or functional consequences, if any, do creators attribute to the session immediately and later, and how do they understand unsuccessful, neutral or ambivalent episodes?
\item[RQ5.] How do musical experience, platform affordances, social context and alternatives to AI shape the meanings attached to these practices?
\end{enumerate}
The questions address participants' accounts of specific events and their interpretations. They do not imply that reported mechanisms are causal or that emotional change necessarily occurs.
